% !PRINTED 95
\weberpage{95}
Verstehen wir also unter $q_1, q_2, q_3 \ldots$ Zahlen, die aus der Reihe der Zahlen 9 und der Primzahlen 7, 13, 19, 31 . . . genommen sind, so ist der allgemeine Ausdruck für $m$:
\weq{1}{m = q_1\,q_2\,q_3\cdots q_\mu.}
Die Indexmoduln sind, wenn $q_1=9$ ist, $c_1=6$, $c_2=q_2-1$, $c_3=q_3-1 \ldots$ Wenn 9 und 7 unter den Factoren von $m$ nicht vorkommen, so fällt der Indexmodul 6 weg.

Die Grössen $\delta_{1,1}, \delta_{1,2}, \ldots \delta_{1,\mu}$ sind die grössten gemeinschaftlichen Theiler von $i_1=3$ mit den $c_1, c_2, \ldots c_\mu$; sie sind also alle $=3$, und es folgt nach \S.\,21, (11):
\[
i_{1,h}=1,\quad c_{1,h}=\tfrac{1}{3}\,c_1.
\]
% [sic: Weber prints c_1 on the right; because h is arbitrary and c_{1,h}=c_h/delta_{1,h}, the intended subscript is h.]
Die Matrix \S.\,21, (24) besteht hier nur aus der einen Verticalreihe $e_{1,1}, e_{1,2}, \ldots e_{1,\mu}$, und diese Grössen kann man beliebig wählen, wenn nur wenigstens eine darunter ist, die durch 3 untheilbar ist. Diese Zahlen $e_{1,1}, e_{1,2}, \ldots e_{1,\mu}$ sind übrigens nur nach dem Modul 3 bestimmt, und können also gleich 0, 1 oder 2 angenommen werden. Die Bedingung 6. des \S.\,21 verlangt aber auch, dass, wenn $q_1=9$ ist, $c_{1,1}e_{1,1}=2e_{1,1}$ nicht durch 3 theilbar sein soll, und wenn $q_2$ eine Primzahl ist, dass
\[
c_{1,2}\,e_{1,2}=\tfrac{1}{3}(q_2-1)\,e_{1,2}
\]
nicht durch $q_2-1$ theilbar sei, d. h. also, es darf \gesperrt{keine der Zahlen} $e_{1,1}, e_{1,2}, \ldots e_{1,\mu}$ durch 3 theilbar sein, und wir können für jede von ihnen $\pm 1$ wählen.

Um also die Indices $\alpha_1, \alpha_2, \ldots \alpha_\mu$ aller Zahlen $a$ zu finden, die in der Periode
\weq{2}{\eta = \sum r^a}
vorkommen können, hat man nach \S.\,22, (3) alle Zahlen $\alpha$ (nach dem Modul 3) zu nehmen, die einer Congruenz
\weq{3}{\pm\alpha_1\pm\alpha_2\pm\alpha_3\cdots\pm\alpha_\mu\equiv 0\ (\mathrm{mod}\ 3)}
genügen, wobei für \gesperrt{eine bestimmte Gruppe eine Vorzeichen-Combination festgehalten werden muss}.

Die Anzahl der möglichen cubischen Körper hängt bei gegebenem $m$ nur von der Zahl $\mu$ der Indexmoduln ab und beträgt, da eines der Vorzeichen in (3) willkürlich angenommen werden kann, $2^{\mu-1}$.

Ganz ähnlich würde sich das Resultat gestalten, wenn man statt 3 eine andere Primzahl $p$ für $i_1$ setzen würde. Es würde dann nur in die Congruenz (3) an Stelle der doppelten Vor-
%CONT>>
